\clearpage
\section{Greitens --- \textit{Surveillance, Security, and Liberal Democracy in the Post-COVID World}}

\subsection{Research question and central argument}

Sheena Chestnut Greitens examines whether the COVID-19 pandemic, and especially the rapid expansion of health surveillance technologies, is likely to transform global patterns of liberalism, democracy, privacy, and civil liberties. The article asks two related questions: whether the pandemic promotes the diffusion of an illiberal model of surveillance associated with China, and whether the use of surveillance itself necessarily produces democratic erosion.

The core argument is deliberately nuanced. China entered the pandemic with a highly developed techno-authoritarian system that integrated surveillance, policing, social control, and state security. COVID-19 strengthened and legitimized this pre-existing architecture. Because China is a major power, a large technology exporter, and an increasingly influential actor in international standard-setting bodies, elements of its model could diffuse internationally. However, diffusion is neither automatic nor uniform. China's governance model depends on institutional features that many other regimes do not possess, which limits wholesale emulation.

More importantly, Greitens argues that surveillance should not be treated as synonymous with autocracy. The pandemic largely intensified pre-existing political trajectories: autocracies were more likely to violate rights, weak democracies and hybrid regimes faced greater risks of backsliding, while many consolidated democracies used intrusive surveillance without significantly undermining democratic institutions. In the latter cases, surveillance was constrained by rule of law, institutional oversight, temporal limits, proportionality, and public accountability. The pandemic's effects on \textit{liberalism} and on \textit{democracy} therefore need to be analytically separated.

\subsection{China's model: prevention, control, and techno-authoritarianism}

China was the first state to confront the outbreak and therefore the first to develop a comprehensive response model. The Chinese concept of \textit{fangkong}, usually translated as ``prevention and control,'' is central to this model. Although the term has public-health applications, Greitens emphasizes that it originated in the language of policing and social stability. Under Xi Jinping, prevention and control became a core principle of the Chinese Communist Party's approach to regime security: threats should be identified, monitored, and neutralized before they generate instability.

The pandemic activated this existing architecture. Public-security officials collaborated with private technology companies such as Alibaba and Tencent to develop health-code systems that combined information on movement, contacts, and health indicators. Citizens were assigned color codes that affected their ability to travel and enter public spaces. These systems were linked to pre-existing government databases and enforced through local officials and neighborhood-level mobilization.

The result was a fusion of public health with public security. Surveillance enabled the state to identify infections and restore mobility, but it also embedded citizens more deeply in a broad, open-ended apparatus of monitoring. Some local governments even proposed retaining health-tracking tools after the acute phase of the pandemic. The response therefore did not create Chinese techno-authoritarianism from scratch; it reinforced and normalized institutions that had already been developed for social control.

Greitens also highlights other illiberal features of the Chinese response. Local officials initially suppressed or failed to transmit information about the outbreak, whistleblowers were silenced, and accountability remained concentrated at lower levels rather than reaching the national leadership. The same conceptual fusion of health and security also has darker precedents: Chinese authorities have used medicalized language to describe political control, including the idea of ``immunizing'' society against ideological or political threats. This illustrates how securitization of health can coexist with the medicalization of repression.

\subsection{Will the Chinese model diffuse?}

The article next uses international-relations theories of diffusion to evaluate whether China's response will spread. Greitens stresses the distinction between \textit{clustering} and \textit{diffusion}. If multiple governments adopt similar policies because they face the same pandemic and possess similar resources, this is clustering. Diffusion requires interdependence: adoption in one country must increase the probability of adoption in another.

China has clear mechanisms through which diffusion could occur. Beijing actively promoted its pandemic response through diplomatic outreach, aid, exports of medical equipment, propaganda, and claims that its governance model had proven effective. Chinese firms already exported surveillance technologies to dozens of countries before COVID-19, creating supplier-client relationships that could facilitate the spread of health-surveillance tools. China's first-mover advantage in testing and refining such technologies also creates market incentives for other governments to purchase them.

International organizations provide another possible channel. Before the pandemic, China had already become more active in global technology standard-setting, including facial-recognition and surveillance standards. Leadership within organizations such as the International Telecommunication Union could help normalize Chinese technical standards and expand market access for Chinese companies. Political and market mechanisms could therefore reinforce each other.

At the same time, diffusion faces important limits. Political practices are most likely to diffuse among countries with comparable regime characteristics, yet China's institutional system is unusually centralized and possesses a dense grassroots apparatus of social control. Many authoritarian regimes lack comparable administrative capacity; democracies face legal and political constraints; federal systems disperse authority; polarized societies may resist uniform implementation; and citizens may tolerate noncompliance to degrees that make Chinese-style enforcement impossible.

The evidence available in the first year of the pandemic also suggested that many of the most visible diffusion processes were not specifically Chinese. Governments learned from Taiwan, South Korea, and other successful democracies. Public-health epistemic communities exchanged information on masking, testing, tracing, and treatment. Thus, successful democratic models can also generate learning and emulation.

\subsection{Pandemic response and democratic erosion}

Greitens broadens the analysis beyond surveillance. The pandemic could damage democracy through numerous mechanisms: emergency powers without sunset clauses, censorship, restrictions on media freedom, postponement or manipulation of elections, abusive enforcement of lockdowns, expanded military involvement in civilian politics, discrimination, and executive aggrandizement. Surveillance is only one possible pathway.

This distinction matters because democratic erosion was already occurring before COVID-19. Leaders such as Viktor Orbán and Rodrigo Duterte had weakened democratic institutions before the pandemic. Analysts must therefore separate trends caused by the pandemic from trends merely accelerated by it or occurring alongside it. The pandemic may act as a permissive condition or provide a convenient justification for further repression without being the original cause of democratic decline.

Early comparative evidence suggests that COVID-19 mainly reinforced existing regime trajectories. The most serious violations of democratic standards occurred disproportionately in states that were already authoritarian or democratically fragile. In contrast, a large majority of consolidated democracies adopted emergency measures without severe violations of democratic standards. The Asian evidence is especially revealing: if simple geographic diffusion from China were decisive, neighboring democracies should have displayed comparable authoritarian responses. Instead, democratic regime type remained strongly associated with lower levels of pandemic-related democratic violations.

\subsection{Democracy-compatible surveillance}

The most important theoretical contribution is the distinction between surveillance and authoritarianism. Consolidated democracies such as Taiwan and South Korea used extensive digital monitoring. Taiwan linked health-insurance data to immigration and customs information and used mobile phones to enforce quarantine. South Korea also relied on smartphone data and warrantless access under legislation revised after the MERS outbreak. Yet both countries avoided major democratic backsliding.

Greitens identifies three characteristics that make surveillance more compatible with democracy. First, measures should be \textit{necessary and proportional}. Surveillance should target actual public-health risks rather than political opponents, and restrictions should be justified by a clear relationship to the health emergency. Second, measures should be \textit{limited in time and scope}. Data retention should be temporary, access restricted to authorized personnel, and emergency systems dismantled when the crisis ends. Taiwan, for example, required deletion of certain data after quarantine periods and promised to remove the monitoring infrastructure after the pandemic. South Korea similarly limited access to integrated data platforms and required eventual deletion of personal information.

Third, surveillance must be subject to \textit{democratic review and accountability}. Courts, legislatures, transparency requirements, and public debate constrain executive discretion. Taiwan's emergency-health framework had been reviewed by the Constitutional Court, while South Korea developed a statutory regime that both empowered and constrained the government. Democratic institutions therefore ``fence in'' surveillance.

Public opinion and political norms are also important. Citizens can accept substantial temporary restrictions when they perceive them as legitimate, effective, and bounded. Transparency can partially offset the informational asymmetry created by surveillance: as the state becomes able to see citizens more clearly, citizens should also be able to scrutinize the state.

\subsection{Implications}

The article concludes that the pandemic does not generate a simple global shift from democracy toward autocracy. China has greater capacity to promote illiberal standards than in the past, but its model is not universally transferable. Democratic systems can employ sophisticated surveillance while preserving accountability if institutions constrain the technology's purpose, duration, access, and political use.

For policy, Greitens recommends focusing democratic assistance on weak democracies where emergency powers are most likely to become permanent. Major democracies and international organizations can also promote democracy-compatible public-health practices, compete in technology standard-setting, and support international norms that protect privacy and civil liberties. The broader research agenda should distinguish carefully among the different mechanisms linking crises, surveillance, liberalism, and democratic erosion.


\clearpage
\section{Hafner-Burton and Tsutsui --- \textit{Human Rights in a Globalizing World: The Paradox of Empty Promises}}

\subsection{The central puzzle}

Emilie Hafner-Burton and Kiyoteru Tsutsui examine a striking contradiction in the international human-rights regime. Since the Second World War, states have ratified an increasing number of international treaties protecting civil, political, economic, social, racial, women's, children's, and physical-integrity rights. Yet the global expansion of formal legal commitments has not been matched by an equivalent decline in government repression. In some periods, the proportion of repressive states even increased while ratification became nearly universal.

This creates a basic puzzle: what is the value of international human-rights treaties if governments can join them while continuing to torture, imprison, disappear, or kill their citizens? The authors argue that neither pessimistic rationalist approaches nor optimistic constructivist and legal approaches fully explain this pattern. Instead, international human-rights institutionalization produces a \textit{paradox of empty promises}.

The paradox has two components. First, weak enforcement makes it easy for governments to ratify human-rights treaties for symbolic legitimacy without changing domestic behavior. Ratification can therefore be ``window dressing,'' creating a sharp decoupling between policy and practice. In some cases, the legitimating benefits of ratification may even coexist with worsening repression. Second, the global institutionalization of human rights simultaneously strengthens human-rights norms and empowers international and domestic civil-society actors. These actors can use the state's own promises as leverage, mobilizing information and reputational pressure that eventually improves behavior.

Thus, formal treaty commitment can be directly ineffective or even negatively associated with rights practices, while the broader human-rights regime indirectly improves behavior through global civil society.

\subsection{Competing theories of compliance}

The article situates this argument within several theoretical traditions. Realist approaches expect little independent impact from international human-rights law because states follow national interests and power considerations. Treaties without credible enforcement should have little capacity to change behavior.

Neoliberal institutionalist approaches are somewhat more optimistic about international regimes but still emphasize rational incentives. Institutions can facilitate cooperation by clarifying obligations, monitoring behavior, and supplying enforcement mechanisms. Human-rights treaties, however, generally lack strong material sanctions, so rationalist theory predicts weak compliance effects.

Liberal approaches emphasize domestic politics. Governments' decisions to commit and comply depend on internal political institutions and interest-group bargaining. Ratification therefore need not guarantee later compliance because domestic coalitions and incentives may change.

Constructivist and international-law perspectives are more optimistic. They emphasize socialization, normative persuasion, and the tendency of states to comply with legal obligations they voluntarily accept. International organizations and nongovernmental actors can redefine legitimate state behavior and socialize elites into new standards.

Hafner-Burton and Tsutsui argue that each perspective captures part of the empirical reality but misses the interaction between formal treaty systems and civil-society activism. Their synthesis draws especially on rational institutionalism and sociological world-society theory.

\subsection{World society, legitimacy, and decoupling}

The world-society approach emphasizes that organizations and states often adopt globally legitimate institutional models because doing so signals modernity and appropriateness, not necessarily because the models are functionally effective. As global standards become taken for granted, organizations converge formally around them. Yet formal adoption can be \textit{decoupled} from actual practice.

Human-rights treaties are a prime example. Once treaty ratification becomes a widely accepted marker of legitimate statehood, nonratifying governments risk appearing deviant. This generates strong incentives to ratify even for governments that lack either the willingness or the capacity to comply. Because enforcement is weak, the symbolic benefits of ratification can be obtained at low cost.

The authors extend ordinary decoupling into what they call \textit{radical decoupling}. The gap between formal commitment and practice can become so severe that ratification is associated with worsening repression. Treaties may provide repressive governments with international legitimacy while failing to constrain violence. A government can point to its legal commitments as evidence of respectability, potentially reducing pressure for deeper reform.

The institutional weakness of the regime is therefore crucial. Human-rights treaties are relatively good at defining norms, collecting information, and establishing reporting procedures, but weak at imposing material or legal sanctions. Their monitoring bodies can expose violations but generally cannot directly punish governments by withholding important resources.

Historical context reinforced this weakness. Many core treaties came into force during the Cold War, when geopolitical alliances often shielded governments from punishment. Ratification could therefore be politically inexpensive, while implementation remained optional.

\subsection{The paradoxical role of global civil society}

The second part of the theory is more optimistic. As more states ratify human-rights treaties, the norms codified in those treaties gain global legitimacy. International nongovernmental organizations (INGOs) such as Amnesty International and Human Rights Watch can then invoke these norms to expose violations, mobilize domestic actors, and create reputational costs for governments.

The key causal mechanism does not run directly from treaty ratification to compliance. Instead, it runs through the growing institutional legitimacy of human rights and the transnational networks that mobilize around them. Governments become increasingly vulnerable to embarrassment and delegitimation when their behavior contradicts internationally accepted standards.

International civil society can also amplify domestic activism. Local organizations can connect to transnational networks, gain information and resources, and publicize abuses abroad. Global actors can then pressure governments in ways that would be difficult for domestic groups acting alone. In effect, civil society supplies some of the enforcement that formal treaty institutions lack.

The authors therefore make two core predictions. First, treaty ratification by itself should have no positive effect on human-rights practices and may even be associated with greater repression. Second, stronger links to international civil society should improve human-rights behavior.

\subsection{Research design and measurement}

The empirical analysis covers 153 states from 1976 to 1999. The dependent variable measures government repression of physical-integrity rights, combining information on murder, torture, forced disappearance, and political imprisonment. The scale ranges from relatively secure rule of law and rare repression to systematic, population-wide terror.

The authors examine ratification of six core UN human-rights treaties: the International Covenant on Civil and Political Rights, the International Covenant on Economic, Social and Cultural Rights, the Convention Against Torture, the Convention on the Elimination of All Forms of Racial Discrimination, the Convention on the Elimination of All Forms of Discrimination against Women, and the Convention on the Rights of the Child. They use both a cumulative count of treaty ratification and treaty-specific measures.

The second key explanatory variable measures the extent to which citizens of a country are connected to international nongovernmental organizations. INGO membership serves as an indicator of embeddedness in global civil society.

The statistical models control for factors already associated with repression, including economic development, trade, democracy, population, previous levels of repression, civil and international conflict, state age, and temporal effects. The use of a long pooled time-series design allows the authors to compare the effects of formal treaty participation and civil-society linkage within the same empirical framework.

\subsection{Main findings}

The findings support the paradox-of-empty-promises argument. Treaty ratification does not reliably improve domestic human-rights practices. Across several specifications, joining more treaties is either unrelated to improved behavior or associated with worsening repression. The results therefore contradict a simple view in which governments become more respectful of rights merely because they formally accept international legal obligations.

This does not mean that the international human-rights regime is irrelevant. The relationship between international civil-society embeddedness and rights practices is consistently more positive. States with stronger connections to INGOs tend to exhibit less repression, even after controlling for domestic political and economic factors.

The most important implication is that the same global institutionalization process has opposite effects through different channels. Legal commitment is easy to make because the treaty regime lacks direct enforcement, which facilitates symbolic ratification and decoupling. Yet widespread ratification also increases the legitimacy and visibility of human-rights norms, which gives nongovernmental actors stronger tools with which to pressure governments.

The authors' Figure 1 captures the empirical puzzle visually: the average percentage of human-rights treaties ratified rises dramatically over time, but the percentage of states classified as repressive does not decline in parallel. This divergence motivates the argument that formal legalization alone is a poor indicator of behavioral change.

\subsection{Theoretical implications}

The article challenges both excessive pessimism and excessive optimism about international human-rights law. Rationalist theories are correct that weak enforcement limits the direct effect of treaties. Governments can make promises they have no intention of keeping. However, rationalist accounts are too pessimistic when they ignore civil society and norm-based reputational pressure.

Constructivist and legal theories are correct that human-rights norms can influence state behavior, but they often attribute too much causal importance to treaty ratification itself. The positive effects of the regime operate more strongly through societal networks and transnational activism than through formal legal obligation alone.

The broader lesson is that international institutions can have indirect and unintended effects. A weak treaty may fail to enforce compliance directly while nevertheless transforming the normative environment in which governments and activists operate. International law can therefore matter even when compliance with its formal rules is poor.

\subsection{Conclusion}

Hafner-Burton and Tsutsui conclude that human-rights globalization is a double-edged process. States frequently make empty promises because ratification brings legitimacy at low enforcement cost. In the short run, formal commitment may therefore be detached from or even negatively associated with practice. But the global spread of those promises also creates a common language of legitimacy that activists can use against governments.

The strongest source of improvement is not treaty ratification in isolation but social embeddedness in a global human-rights community. Effective human-rights governance therefore depends on strengthening the connections between international norms, civil society, information networks, and domestic political mobilization rather than assuming that formal legal commitments automatically produce compliance.


\clearpage
\section{Murdie and Davis --- \textit{Shaming and Blaming: Using Events Data to Assess the Impact of Human Rights INGOs}}

\subsection{Research question and argument}

Amanda Murdie and David Davis investigate whether ``naming and shaming'' by human-rights international nongovernmental organizations (HROs) actually improves state human-rights practices. Organizations such as Amnesty International and Human Rights Watch routinely publicize abuses and criticize governments, and such activity occupies a central place in theories of transnational advocacy. Yet global quantitative evidence of its effectiveness had been limited.

The authors' main argument is conditional: \textit{shaming by itself is not enough}. HRO criticism becomes effective when it is embedded in a broader advocacy network that generates pressure from both below and above. Pressure ``from below'' comes from the domestic presence of human-rights organizations that can support local mobilization. Pressure ``from above'' comes from third-party states, intergovernmental organizations, individuals, and other actors that take up HRO claims and pressure the targeted government.

The article therefore provides a quantitative test of the boomerang model associated with Keck and Sikkink and the spiral model associated with Risse, Ropp, and Sikkink. Both frameworks argue that international advocacy is most effective when domestic and external forces interact.

\subsection{How HROs influence governments}

The theoretical literature assigns HROs several functions. First, they place human-rights violations on the international agenda. Second, they provide information, resources, legal assistance, and organizational support to domestic activists. Third, they encourage other governments and international organizations to join pressure campaigns.

Through shaming, HROs seek to expose the gap between a government's claimed identity as a legitimate member of international society and its actual behavior. This creates reputational pressure. Governments may initially make tactical concessions simply to reduce criticism, but repeated interaction can produce deeper socialization and eventually ``rule-consistent behavior.''

The boomerang model describes the process by which domestic groups that are blocked from influencing their government directly reach outward to transnational allies. International HROs publicize their claims and mobilize external actors, who then pressure the government from outside. The resulting external pressure reinforces domestic mobilization.

The spiral model adds a temporal dimension. Governments may initially deny abuses and invoke sovereignty, then make cosmetic concessions, later accept the validity of international norms, and eventually institutionalize compliance. In both models, isolated shaming is insufficient; it must be connected to domestic organization and broader international pressure.

\subsection{Hypotheses}

The authors derive four hypotheses. First, HRO shaming alone should not systematically improve human-rights performance. Second, shaming should become more effective when many HROs are present inside the target state. Third, shaming should become more effective when third-party actors also target the government. Fourth, the strongest improvements should occur when shaming is combined with both domestic HRO presence and third-party targeting.

These hypotheses distinguish the article from studies that simply count the number of NGOs in a country or examine the output of a single organization. The authors seek to measure what HROs actually \textit{do} and how that activity interacts with the broader advocacy environment.

\subsection{Events data and research design}

The study constructs a new dataset covering 1992--2004. It identifies more than 400 human-rights organizations from the \textit{Yearbook of International Organizations} and uses the Integrated Data for Event Analysis (IDEA) system to code Reuters reports involving these organizations.

Events are coded in a ``who did what to whom'' structure. The authors identify conflictual actions directed by HROs against governments, such as public criticism or demands. Aggregated to the country-year level, this becomes the measure of HRO Shaming. The dataset contains thousands of HRO-government interactions and provides far broader coverage than measures based only on Amnesty International.

The authors also construct a measure of \textit{Indirect Targeting}, which captures cases where a third-party actor criticizes a government while citing information or claims generated by an HRO. Such third parties include states, intergovernmental organizations, prominent individuals, and other organizations. This measure operationalizes pressure from above.

Domestic HRO presence is measured by the number of human-rights organizations with members or volunteers in the target country. The dependent variable is the Cingranelli--Richards Physical Integrity Rights Index, which measures respect for freedom from torture, extrajudicial killing, political imprisonment, and disappearance. The authors analyze both the level of human-rights performance and whether a country experiences an improvement from one year to the next.

The statistical models control for economic development, democracy, population, and war. Robustness tests address autocorrelation, media-coverage bias, alternative measures of repression, endogeneity concerns, and different temporal lags.

\subsection{Results: shaming alone does not work}

The first finding strongly supports the conditional argument. HRO shaming by itself has no statistically reliable positive effect on subsequent human-rights performance. A government can be criticized repeatedly without changing its behavior if the criticism is not accompanied by domestic organizational capacity or additional pressure from influential third parties.

This finding helps reconcile apparently pessimistic earlier studies with transnational-advocacy theory. If theory never predicted an unconditional shaming effect, then evidence showing that isolated criticism fails is not necessarily evidence against the boomerang or spiral models. The correct empirical test must examine interactions.

\subsection{Domestic pressure from below}

The interaction between HRO shaming and domestic HRO presence is positive. Shaming becomes significantly more effective when a substantial number of human-rights organizations operate in the target state. The authors estimate that a relatively high domestic HRO presence---roughly around fifty organizations in one interpretation of the model---is sufficient for shaming to begin producing meaningful improvements.

The mechanism is consistent with pressure from below. Domestic HROs connect international criticism to local mobilization, provide resources and information, assist victims, and make it harder for governments simply to ignore external accusations. International publicity gains political leverage when actors on the ground can translate it into sustained pressure.

\subsection{Third-party pressure from above}

The interaction between HRO shaming and indirect targeting also improves human-rights outcomes. When states, international organizations, or prominent individuals echo HRO accusations, the reputational and political costs of repression increase. Third-party involvement transforms a report by an NGO into a broader international campaign.

This result supports one of the most important claims of transnational-advocacy theory: HROs matter not only because they directly criticize governments but because they recruit other actors into the campaign. Their informational and agenda-setting role changes the behavior of actors with greater diplomatic, institutional, or political leverage.

The combination of domestic presence and third-party pressure produces the strongest effects. As the overall advocacy environment becomes denser, the marginal effect of HRO shaming becomes increasingly positive. The article therefore finds substantial support for the ``from above and below'' logic.

\subsection{Economic vulnerability and reputation}

The authors also test a competing explanation in which shaming works mainly when governments depend heavily on foreign aid, trade, or investment and therefore fear material punishment. Their results do not show that economic vulnerability is the central condition making HRO shaming effective.

Instead, the findings point to reputational and political mechanisms. Governments can be vulnerable to criticism because they care about legitimacy and international standing, even when they are not heavily dependent on external economic resources. This strengthens the constructivist dimension of the boomerang and spiral models while still leaving room for material pressure in specific cases.

\subsection{Contribution and conclusion}

The article makes both methodological and theoretical contributions. Methodologically, it demonstrates the value of events data for studying nonstate actors. Rather than relying on organizational presence or the output of one famous NGO, the dataset records actual public shaming behavior by hundreds of organizations and links this activity to third-party reactions.

Theoretically, the study shows that advocacy networks, rather than isolated organizations, are the appropriate unit for understanding human-rights change. HROs are effective when they are embedded in broader domestic and international coalitions. Shaming generates information and reputational pressure, but these mechanisms require local organizational capacity and/or external amplification.

The central conclusion is therefore clear: \textit{naming and shaming works conditionally}. Public criticism alone is weak, but when HROs can mobilize domestic actors and recruit third parties, their campaigns can produce measurable improvements in physical-integrity rights.


\clearpage
\section{von Stein --- \textit{Making Promises, Keeping Promises: Democracy, Ratification and Compliance in International Human Rights Law}}

\subsection{The puzzle of commitment and compliance}

Jana von Stein asks when international human-rights agreements (HRAs) actually change behavior. The article begins from a familiar puzzle: human-rights treaties usually lack strong international enforcement mechanisms, yet some governments comply with them while others treat ratification as an empty promise.

The author argues that the literature has often conflated two distinct meanings of democracy. First, democracy is a set of rights and practices: free elections, civil liberties, due process, and related protections are themselves components of many human-rights treaties. Second, democratic institutions are mechanisms of \textit{enforcement}: independent courts, civil society, and elections make it costly for leaders to violate commitments.

When scholars study treaties such as the International Covenant on Civil and Political Rights, these two dimensions overlap. Democracies already tend to respect many of the rights in the treaty, so it is difficult to determine whether ratification causes better behavior or merely reflects pre-existing democratic practices. Von Stein therefore proposes studying rights that are not definitional to democracy. Her empirical case is child labor and the International Labour Organization's Minimum Age Convention (MAC).

The core argument is two-sided. Domestic enforcement makes treaty commitments potentially powerful, but that same prospect of enforcement can deter governments with poor rights practices from ratifying in the first place. Democracies that expect difficulty complying may avoid the treaty because breaking the promise would be politically and legally costly. Nevertheless, when democracies with substantial child labor do become parties, ratification can produce real improvements. In nondemocracies, by contrast, ratification is usually easier because domestic enforcement mechanisms are weak, but the promise is less likely to be kept.

\subsection{Compliance versus enforcement}

Von Stein separates \textit{compliance} from \textit{enforcement}. Compliance refers to the degree to which domestic practice matches treaty obligations. Enforcement refers to the mechanisms that impose costs on governments when they fail to comply.

Three domestic enforcement mechanisms receive particular attention. First, independent courts allow citizens and groups to challenge government behavior and can make continued violation legally costly. Even when court rulings do not generate perfect compliance, they increase the difficulty of ignoring legal obligations.

Second, civil society can monitor, publicize, litigate, and mobilize around treaty commitments. Ratification can strengthen activists because it makes it harder for governments to deny that they have accepted a principle. It may also create new legal opportunities and international allies.

Third, democratic elections create political accountability. Citizens may punish governments for appearing incompetent, unlawful, or untrustworthy. Treaty compliance rarely decides elections by itself, but breaking an international commitment can become a valence issue that damages a government's reputation for competence and respect for law.

These mechanisms mean that democracies face real domestic costs when they make promises they cannot keep. The same enforcement capacity that creates compliance after ratification can therefore reduce the probability of ratification before the commitment is made.

\subsection{Why study child labor?}

Child labor is analytically useful because it is not a defining feature of democracy. A country can have competitive elections, civil liberties, and independent courts while still having significant child labor. This allows the author to distinguish the institutional conditions that enforce a treaty from the substantive right protected by that treaty.

The Minimum Age Convention was adopted in 1973 and seeks the effective abolition of child labor. Unlike some human-rights agreements, it is relatively precise, does not permit reservations, and covers an outcome for which cross-national data exist. More than 85 percent of states eventually became parties, but ratification timing varied widely, and some countries with extensive child labor delayed or avoided participation.

\subsection{Ratification logic}

The author expects democracies to weigh compliance costs more heavily than nondemocracies. If a democratic country has high levels of child labor, ratification creates a potentially enforceable promise that will be difficult to fulfill. Leaders therefore have incentives to delay or avoid commitment.

However, ratification decisions are multidimensional. Governments face normative pressure from NGOs, regional peers, and international organizations. States embedded in the broader children's-rights or ILO treaty system may be more likely to ratify. Governments may also support the treaty because they identify with its goals or wish to signal commitment to international standards. Political circumstances and institutions can change over time, so some countries may enter the treaty under conditions that later make enforcement more consequential.

Nondemocracies face a different calculus. Courts, civil society, and electoral accountability are weaker, so a commitment is less likely to generate domestic costs. Repressive or noncompliant regimes may therefore ratify without intending to implement the treaty. For them, treaty membership can be relatively cheap symbolic behavior.

\subsection{Method and variables}

The empirical analysis examines both ratification and compliance. Ratification timing is modeled with survival analysis. The key explanatory variables are child-labor levels and three measures of domestic enforcement: judicial independence, civil liberties, and democratic elections. The author interacts child labor with each enforcement measure to test whether poor substantive performance deters ratification more strongly when enforcement is credible.

The ratification models also include international embeddedness, regional ratification, common-law legal heritage, and domestic ratification barriers. These variables capture normative, peer, legal, and institutional influences that might affect treaty commitment.

The compliance analysis examines whether MAC ratification reduces child labor. Crucially, ratification is endogenous: states choose whether and when to join. If democracies with low child labor are the states most likely to ratify, a simple comparison would exaggerate treaty effectiveness. Von Stein therefore uses instrumental-variable methods to separate the causal effect of ratification from selection into the treaty.

The models control for economic development, state capacity, trade exposure, war, temporal trends, and other factors known to affect child labor. This is important because poverty and weak administrative capacity may make compliance difficult even in democratic systems.

\subsection{Ratification findings}

The results show that enforcement concerns shape democratic ratification. Among countries with strong civil liberties or democratic elections, higher child labor significantly reduces the probability of ratifying the Minimum Age Convention. A one-unit increase in child labor makes countries with the strongest civil-liberties protections substantially less likely to ratify; a similar pattern appears for democratic elections.

This is exactly what the enforcement argument predicts. When leaders expect activists, voters, and institutions to hold them accountable for broken promises, they become more cautious about joining agreements that require substantial policy change.

The evidence for judicial independence is weaker. The estimated relationship generally points in the expected direction, but it is not consistently statistically significant. Von Stein discusses possible explanations, including measurement limitations and the possibility that access to courts rather than judicial independence alone is the relevant mechanism.

The broader ratification results also reveal international social pressures. Countries already embedded in children's-rights and ILO institutions ratify more quickly. Common-law countries tend to ratify more slowly. Interestingly, democracies with little or no child labor are especially likely to ratify, consistent with the idea of norm exportation: states that already satisfy the treaty can support it at relatively low cost and promote their preferred standards abroad.

\subsection{Compliance findings}

The compliance results are the article's most important contribution. Once selection into ratification is taken seriously, participation in the Minimum Age Convention is associated with significant reductions in child labor in countries possessing robust democratic enforcement institutions.

The effect of ratification depends on the strength of judicial independence, civil liberties, and democratic elections. Where these institutions are strong, treaty membership produces meaningful improvement. Where they are weak, the estimated effect approaches zero. The article's Figure 2 illustrates this conditional relationship: the predicted reduction in child labor after ratification becomes more substantial as democratic enforcement mechanisms strengthen.

This means that international law can do more than simply reflect pre-existing preferences. Even democracies that initially have serious child-labor problems can improve after taking on a binding international commitment, provided domestic actors possess the tools necessary to enforce the promise.

For nondemocracies, the findings are much less optimistic. Ratification does not reliably reduce child labor. Without independent courts, protected civil society, or electoral accountability, citizens have few mechanisms for transforming international commitments into domestic political costs. Treaty ratification in these cases resembles the ``empty promise'' described in earlier human-rights research.

Economic development and state capacity also play important roles. Poor or administratively weak states can struggle to eliminate child labor even when they are politically accountable. This reinforces the need to distinguish enforcement from capacity: willingness and political pressure are not sufficient when governments lack the bureaucratic and financial means to implement policy.

\subsection{Broader theoretical implications}

The article highlights a selection problem in international institutions. The treaties most capable of changing behavior may deter precisely those governments for which compliance would require the largest reforms. If domestic enforcement is strong, leaders with serious rights problems may choose not to ratify. Looking only at parties can therefore produce misleading conclusions about treaty effectiveness.

At the same time, selection does not make international law irrelevant. Governments face multiple pressures and may ratify despite expected compliance costs; institutions change; and countries can democratize after ratification. Once a meaningful commitment exists in a setting with domestic enforcement, civil society, courts, and elections can use the treaty to produce real behavioral change.

The results also clarify why previous studies of human-rights treaties have produced mixed findings. When the substantive rights protected by a treaty overlap heavily with the institutions that enforce it, it is difficult to separate treaty effects from democratic selection. Studying rights such as child labor reveals the causal role of domestic enforcement more clearly.

\subsection{Conclusion}

Von Stein concludes that international human-rights agreements can matter, but their effects are conditional on domestic institutions. Democracies are more cautious about making costly promises because those promises are enforceable. Yet when they do ratify, strong accountability mechanisms can convert international law into meaningful improvements. Nondemocracies can make the same promises at lower political cost, but the absence of domestic enforcement means they are seldom kept.

The larger lesson for international relations is that commitment and compliance must be analyzed together. Whether international law changes behavior depends not only on treaty design or state preferences but also on the domestic institutions that determine who can invoke the law, impose costs for violation, and transform formal commitments into actual policy.
